\Needspace{16\baselineskip}
\FormulaThmDeltaK[Die Potenzen im Wiederholungsprädikat sind typisiert]{%
 \begin{aligned}[t]
 &\SemigroupAlg{A}{\star}\dsep a\in A\dsep i,j\in\mathbb N
 \dsep i\neq0\dsep j\neq0\vdash{}\\[-2pt]
 &a^i\in A\land a^j\in A
 \end{aligned}
}{B37IPPositivePowersTyped}{\DeltaRow{Mengen}{A\dsep a\dsep i\dsep j}\DeltaRow{Funktionen}{\star}}
\begin{tabproofwide}
 \proofstepwidestar[1]{\SemigroupAlg{A}{\star}}{\rA}
 \proofstepwidestar[2]{a\in A}{\rA}
 \proofstepwidestar[3]{i,j\in\mathbb N}{\rA}
 \proofstepwidestar[4]{i\neq0}{\rA}
 \proofstepwidestar[5]{j\neq0}{\rA}
 \proofstepwidestar[1,2,3,4]{a^i\in A}{\FormulaRefAuto{SemigroupPowerClosure}[thm]{1,2,\rAEa{3},4}}
 \proofstepwidestar[1,2,3,5]{a^j\in A}{\FormulaRefAuto{SemigroupPowerClosure}[thm]{1,2,\rAEb{3},5}}
 \proofstepwidestar[1,2,3,4,5]{a^i\in A\land a^j\in A}{\rAI{6,7}}
\end{tabproofwide}
Im Wiederholungsprädikat werden beide Potenzen unter genau diesen
Typ- und Positivitätsbedingungen verglichen. Eine Potenz mit Exponent null
ist somit für die Definition nicht erforderlich.

\Needspace{18\baselineskip}
\FormulaThmDeltaK[Die Wiederholungsmenge ist eindeutig bestimmt]{%
 \begin{aligned}[t]
 &\SemigroupAlg{A}{\star}\dsep a\in A\vdash{}\\[-2pt]
 &\Bigl(\exists!R\,\forall j\,\Bigl(j\in R\leftrightarrow
     \bigl(j\in\mathbb N\land j\neq0\land{}\\[-2pt]
 &\qquad\exists i\in\mathbb N\,
       (i\neq0\land i<j\land a^i=a^j)\bigr)\Bigr)\Bigr)
       \land\mathcal R_{A,\star}(a)\subseteq\mathbb N
 \end{aligned}
}{B37IPRepetitionSetWellDefined}{\DeltaRow{Mengen}{A\dsep a\dsep i\dsep j\dsep R}\DeltaRow{Funktionen}{\star}}
In dieser Tabelle kürzen wir
\[
 Q(j)\equiv j\neq0\land\exists i\in\mathbb N\,
   (i\neq0\land i<j\land a^i=a^j)
\]
ab. Die eindeutige Existenz folgt aus Aussonderung und Extensionalität;
sie benötigt noch keine Endlichkeit von \(A\).
\begin{tabproofwide}
 \proofstepwidestar[1]{\SemigroupAlg{A}{\star}}{\rA}
 \proofstepwidestar[2]{a\in A}{\rA}
 \proofstepwidestar[]{\exists R\,\forall j\,(j\in R\leftrightarrow j\in\mathbb N\land Q(j))}{\FormulaRefAuto{\exists B\;\forall x\;\bigl(x \in B \;\leftrightarrow\; x \in A \,\land\, P(x)\bigr)}[ax]}
 \proofstepwidestar[]{\exists!R\,\forall j\,(j\in R\leftrightarrow j\in\mathbb N\land Q(j))}{\FormulaRefAuto{\exists A(\forall x(x \in A\leftrightarrow P(x)))\vdash \exists! A(\forall x(x \in A\leftrightarrow P(x)))}[thm]{3}}
 \proofstepwidestar[5]{j\in\mathcal R_{A,\star}(a)}{\rA}
 \proofstepwidestar[1,2,5]{j\in\{v\in\mathbb N\mid Q(v)\}}{\FormulaRefAuto{FiniteSemigroupPowerRepetitionSetDef}[def]{1,2,5}}
 \proofstepwidestar[1,2,5]{j\in\mathbb N\land Q(j)}{\FormulaRefAuto{x\in\{u\in A\mid P(u)\}\eqvdash x\in A\land P(x)}[thm]{6}}
 \proofstepwidestar[1,2,5]{j\in\mathbb N}{\rAEa{7}}
 \proofstepwidestar[1,2]{\mathcal R_{A,\star}(a)\subseteq\mathbb N}{\FormulaRefAuto{A\subseteq B\coloneq\forall x\,(x\in A\to x\in B)}[def]{\rUI{\rRI{5,8}}}}
 \proofstepwidestar[1,2]{\begin{aligned}[t]&\bigl(\exists!R\,\forall j\,(j\in R\leftrightarrow j\in\mathbb N\land Q(j))\bigr)\\[-2pt]&{}\land\mathcal R_{A,\star}(a)\subseteq\mathbb N\end{aligned}}{\rAI{4,9}}
\end{tabproofwide}
